\chapter[Introdução]{Introdução}

O \textit{Micromouse} é uma competição de robótica móvel autônoma na qual um pequeno robô deve percorrer um labirinto desconhecido e alcançar uma região-alvo, geralmente localizada no centro do labirinto, no menor tempo possível. A proposta surgiu no contexto da competição \textit{Amazing MicroMouse Maze Contest}, que foi idealizada pela \textit{IEEE Spectrum}, na qual os robôs deveriam possuir lógica e memória próprias para explorar o labirinto, aprender com tentativas anteriores e melhorar o desempenho em uma corrida final \cite{christiansen2014amazing}. Desde então, o desafio tornou-se uma atividade recorrente em eventos acadêmicos e competições de robótica, passando a ser como plataforma para o estudo de navegação autônoma, planejamento de trajetórias, sensoriamento, controle de motores e integração de sistemas embarcados \cite{milos2010undergraduate}.

De acordo com regras tradicionais da modalidade, o labirinto padrão do \textit{Micromouse} é composto por uma matriz de células, comumente organizada em uma grade de $16 \times 16$, em que cada célula possui dimensões de $18 \times 18$ cm. As paredes possuem aproximadamente 5 cm de altura e 1,2 cm de espessura, enquanto o robô não deve exceder 25 cm de comprimento ou largura \cite{ieeerhu2024micromouse}. Além disso, o robô deve operar sem intervenção humana durante a execução, ou seja, não sendo possível receber informações externas sobre a configuração do labirinto após o início da prova. No contexto deste projeto, proposto pela disciplina de PI1, o \textit{micromouse} deverá alcançar zonas em três labirintos desconhecidos, com dimensões de $4 \times 4$, $8 \times 8$ e $16 \times 16$ células, respeitando as restrições físicas e operacionais propostas para a competição \cite{garcia2026tema}.

O desenvolvimento de um \textit{Micromouse} envolve a integração de diferentes áreas da engenharia. Milos et al. \cite{milos2010undergraduate} descrevem esse tipo de projeto como uma atividade interdisciplinar, na qual conhecimentos de diferentes áreass, como ciência da computação, engenharia de computação, engenharia elétrica e engenharia mecânica são aplicados conjuntamente para transformar conceitos teóricos em um sistema físico funcional. Essa característica também está presente no presente projeto, que demanda a concepção da estrutura mecânica do robô, a seleção de sensores e atuadores, o desenvolvimento do circuito eletrônico, a programação embarcada, o controle de movimento e a implementação de um sistema de visualização e armazenamento dos dados obtidos durante a execução, abrangendo todos os cursos de Engenharia da FCTE.

Assim, o desenvolvimento do \textit{Micromouse} é justificado, nesse contexto, como uma solução educacional e tecnológica capaz de reunir, em um único produto, conceitos de robótica móvel, sistemas embarcados, eletrônica, controle, programação, banco de dados e desenvolvimento web. O produto proposto não se limita à construção de um robô que percorre labirintos, pois também inclui a exibição em tempo real de dados de telemetria, como tipo de labirinto, trajeto percorrido, consumo de bateria, velocidade média, tempo de conclusão e indicação de cumprimento do desafio, além do armazenamento dos resultados finais para consultas posteriores \cite{garcia2026tema}. Com isso, o projeto atende às necessidades da disciplina ao propor um sistema próprio, verificável e interdisciplinar, permitindo avaliar tanto o desempenho físico do robô quanto a qualidade das informações geradas durante sua operação.

O trabalho presente neste documento está estuturado da seguinte forma: o capítulo \ref{termo} apresenta o termo de abertura, com informações importantes para o projeto; a equipe de trabalho, e a avaliação de desempenho dos memebros, é apresentada no capítulo \ref{equipe}; o capítulo \ref{projeto} apresenta o projeto conceitual do produto, que define as características das quatros áreas presentes no projeto (Estrutura, Energia, Eletrônica e Software), bem como a integração entre elas. No capítulo \ref{orca}, é apresentado o orçamento do projeto; e o cronograma é documentado no capítulo \ref{crono}. Por fim, o capítulo \ref{conclusao} apresenta a síntese do documento.